%% generated by results/harvest.py -- do not edit
\begin{figure}
\centering
\begin{tikzpicture}[x=1cm, y=1cm]
\node[art]  (m)  at (2.4, 0)    {55 applicable mutants of the
  \texttt{riscv-btor2} emissions};
\node[tool] (g1) at (2.4, -1.0) {square suite (conjoined probes)};
\node[tool] (g2) at (2.4, -2.0) {authored branch questions};
\node[tool] (g3) at (2.4, -3.0) {derived benchmark (78 questions)};
\node[art]  (e)  at (2.4, -4.0) {\textbf{escaped: 0}};
\node[art] (k1) at (6.2, -1.0) {51 killed};
\node[art] (k2) at (6.2, -2.0) {0 killed};
\node[art] (k3) at (6.2, -3.0) {4 killed};
\draw[chk] (m) -- (g1);
\draw[chk] (g1) -- node[elbl, right] {4 survive} (g2);
\draw[chk] (g2) -- node[elbl, right] {4 survive} (g3);
\draw[chk] (g3) -- (e);
\draw[chk] (g1) -- (k1);
\draw[chk] (g2) -- (k2);
\draw[chk] (g3) -- (k3);
\end{tikzpicture}
\caption{The escape-rate experiment (\Cref{tab:escape}) as a gate
stack: every applicable mutant runs the gates in the order the
platform applies them, and each gate's kill count is measured. The
authored branch questions catch nothing the square suite missed ---
the derived benchmark is the differentiator (a finding in itself).}
\label{fig:escape}
\end{figure}
